\documentclass{bht-thesis}

\usepackage[ngerman]{babel}
\usepackage[utf8]{inputenc} % Required for inputting international characters
\usepackage[T1]{fontenc} % Output font encoding for international characters
\usepackage{mathpazo} % Use the Palatino font by default
\usepackage[scaled]{helvet}
\usepackage{eulervm}

\renewcommand\bf{\bfseries}
\usepackage[longnamesfirst]{natbib}
\setcitestyle{citesep={;}, aysep={,}}
\bibliographystyle{agsm}

\usepackage{blindtext}

% Standard Page settings
\geometry{
	% showframe,
	left=2.5cm,
	right=2.5cm,
	top=2.5cm,
	bottom=2.5cm
	}

% Information about the thesis
\title{Master Thesis}
\degree{Master of Science}
\university{Beuth Hochschule für Technik Berlin}
\faculty{Fachbereich VI - Informatik und Medien}
\course{Medieninformatik}

\thesistitle{Titel der Abschlussarbeit}
\author{Peter Pan}

\address{Bla Str. 123\\14852 Berlin}
\matnr{123456}
\phone{+49 1234 56789}
\email{peter@pan.de}
\city{Berlin}

\supervisor{Prof. Bla}
\expert{Prof. Blubb}

\date{24. Dezember 2019}

\newcounter{savepage}

% Document starts here
\begin{document}

% Titlepage and listings
\cleardoublepage
\pagenumbering{Roman}
\maketitle
\tableofcontents
\listoffigures
\listoftables
\lstlistoflistings
\cleardoublepage

% Real document starts here
\onehalfspacing
% Save roman page number in savepage
\setcounter{savepage}{\arabic{page}}
\pagenumbering{arabic}

\chapter{Einleitung}
\section{Jo}
\blindtext[4]

\blindtext[2]
\begin{table}
	\centering
	\begin{tabular}{|l|l|}
		\hline
		Hallo & Hallo \\
		\hline
		Jo & Jo \\
		\hline
	\end{tabular}
	\caption{This is my one big table} \label{tab:sometab}
	\end{table}
\blindtext[3]
\section{Jojo}
\blindtext[2]

\blindtext[3]\footnote{Test einer Fußnote}

\blindtext[5]

\blindtext[1]\footnote{Test einer Fußnote}

\section{Lala}
\blindtext[3]
\begin{lstlisting}[language=Java,float,caption={Nice Java class},label={lst:class}]
	public class MyClass{
		public MyClass(){
		// creates a new instance
		}
	}
\end{lstlisting}
\blindtext[4]
\blindtext[2]
\chapter{Fazit}
\section{Bla}
\blindtext[1]
\begin{figure}
	\includegraphics{figures/Beuth-Logo_basis}
	\centering
	\caption{Beuth Logo}
	\label{fig:bht-logo}
\end{figure}
\blindtext[3]
\blindtext[5]
\subsection{Blablubb}
\blindtext[4]
\blindtext[2]
\blindtext[3]
\subsection{Blubbobo}
Wie \citet[S.~6]{ahu61} bereits erläuterte.

Und nochmal wurde es so gemacht \citep[S.~15]{ahu61}.

Es geht ab \cite[S.~121]{ahu61}.

Wie \citet[S.~6]{ahu61} bereits erläuterte.

Und nochmal wurde es so gemacht \citep[S.~15]{ahu61}.

Es geht ab \cite[S.~121]{ahu61}.

% Main document end
\cleardoublepage
% Switch back to roman numbering
\pagenumbering{Roman}
% Read pagenumber from previously saved value savepage
\setcounter{page}{\thesavepage}

% Bibliography
\bibliography{main}
% \printbibliography

% Appendix
\cleardoublepage
\appendix
\chapter{Anhang}
\section{Wichtiger Anhang}
Ganz wichtiger Anhang.
\section{Noch ein wichtiger Anhang}
Super wichtiger Anhang.

% Statuary declaration
\cleardoublepage
\pagenumbering{gobble}
\statutorydeclaration
\end{document}